\section{Conclusion}
\label{sec:conclusion}

We evaluate whether simple, handcrafted adversarial prompts placed after sensitive information in LLM system prompts can reduce information leakage.
Across 5{,}650 API calls on \gptmini with 65 diverse extraction attacks, we find that our best design---a structured boundary marker (\modfirewall)---reduces exact sensitive field leakage by 58\% (from 15.3\% to 6.4\%, $p < 10^{-7}$), matching the effectiveness of standard heuristic guardrails without any optimization or training.
Two of three firewall designs achieve statistically significant protection.

Contrary to our initial hypothesis, firewalls do not create selective permeability: they broadly suppress leakage for both exact-match and fuzzy queries rather than allowing query-type-dependent access control.
We also find that injection-style attacks account for nearly all successful extractions on \gptmini, while polite-request attacks produce near-zero leakage even without defenses.

Our key takeaway is that even zero-cost, handcrafted adversarial firewalls provide meaningful security improvements for deployed LLM applications, offering a practical first line of defense that can be combined with more sophisticated protections.

Future work should extend this evaluation to multiple models (Claude, Llama-3, Gemini), test against adaptive adversaries that specifically target the firewall mechanism, evaluate multi-turn attack escalation, and explore LLM-optimized firewalls that build on the handcrafted designs as starting points.
Mechanistic analysis on open-source models---examining how firewall text affects attention patterns to sensitive tokens---could provide deeper understanding of why these simple defenses work and how to make them more robust.
